% (0,1) 的开覆盖 U_n=(1/n,1)，n=2,3,4,5,8；有限子族的并 (1/N,1) 在 0 附近留下缺口 (0,1/N]，此处 N=8。数轴单位 8cm。
\begin{tikzpicture}[x=8cm,
  note/.style={font=\FigureNoteFont\fontsize{7.5}{9}\selectfont,text=FigureMuted,align=center}]
  \draw[FigureStroke,line width=.6pt] (-.03,0)--(1.05,0);
  \draw[FigureStroke,fill=white,line width=.8pt] (0,0) circle (2pt);
  \draw[FigureStroke,fill=white,line width=.8pt] (1,0) circle (2pt);
  \node[below=4pt] at (0,0) {$0$}; \node[below=4pt] at (1,0) {$1$};
  \foreach \n [count=\k] in {2,3,4,5,8}{
    \pgfmathsetmacro{\a}{1/\n}
    \pgfmathsetmacro{\yy}{0.32*\k}
    \draw[FigureBlue,line width=1.6pt] (\a,\yy)--(1,\yy);
    \draw[FigureBlue,fill=white,line width=.7pt] (\a,\yy) circle (1.6pt);
    \draw[FigureBlue,fill=white,line width=.7pt] (1,\yy) circle (1.6pt);
    \node[anchor=west,font=\footnotesize] at (1.015,\yy) {$U_{\n}$};
    \draw[FigureGrid,line width=.4pt,densely dotted] (\a,0)--(\a,\yy);
  }
  \draw[FigureGreen,line width=2.4pt] (0.004,-.02)--(.125,-.02);
  \node[note,text=FigureGreen!80!FigureInk,anchor=west] at (.18,-.2) {缺口$(0,\tfrac18]$};
  \node[below=4pt,font=\footnotesize] at (.125,0) {$\tfrac18$};
  \node[note,anchor=west] at (.3,1.92) {左端点$1/n$逼近$0$，但每个$U_n$都不含$0$附近的全部点};
\end{tikzpicture}
